% subsection content: 等比型 \, $a_{n+1}=ra_n$

次に説明するのは等比型である.
一般的な解法も先ほどと同じで,\,$a_{n+1}$から順に$a_1$に向けて書き下して行くことで求めることができる.
\begin{align*}
   & a_{n+1}=ra_n                             \\
   & a_{n+1}=r(ra_{n-1})                      \\
   & a_{n+1}=r\{r(ra_{n-2})\}                 \\
   & \qquad \vdots                            \\
   & a_{n+1}=r\cdot r \cdot \cdots \cdot ra_1 \\
   & \therefore \quad a_n=r^{n-1}a_1
\end{align*}

\noindent \textbf{【例題】}
\begin{enumerate}[label=(\arabic*),parsep=3pt]
  \item $\displaystyle a_1=1,\, a_{n+1}=2a_n$
  \item $\displaystyle a_1=1,\, a_{n+1}=-a_n$
  \item $\displaystyle a_1=2,\, a_{n+1}=\frac{1}{2}a_n$
\end{enumerate}

\noindent\textbf{【方針】}
\begin{enumerate}[label=(\arabic*),parsep=3pt]
  \item 等比型の一般項
        $\displaystyle a_n=a_1r^{n-1}$
        を用います.
  \item 公比は$r=-1$です.
        符号が交互に変化することに注意します.
  \item 初項と公比を一般項に代入します.
\end{enumerate}

\noindent\textbf{【解答】}
\begin{enumerate}[label=(\arabic*),parsep=3pt]
  \item $\displaystyle a_n=2^{n-1}$
  \item $\displaystyle a_n=(-1)^{n-1}$
  \item $\displaystyle a_n=2\left(\frac12\right)^{n-1}$
\end{enumerate}

\noindent\textbf{【解法】}\\
解法略.
